
\documentclass{article}
\usepackage{colortbl}
\usepackage{makecell}
\usepackage{multirow}
\usepackage{supertabular}

\begin{document}

\newcounter{utterance}

\centering \large Interaction Transcript for game `cladder', experiment `full\_v1.5\_default', episode 5585 with score{-}Qwen3.5{-}9B{-}sp.
\vspace{24pt}

{ \footnotesize  \setcounter{utterance}{1}
\setlength{\tabcolsep}{0pt}
\begin{supertabular}{c@{$\;$}|p{.15\linewidth}@{}p{.15\linewidth}p{.15\linewidth}p{.15\linewidth}p{.15\linewidth}p{.15\linewidth}}
    \# & \multicolumn{2}{c}{Player} && \multicolumn{2}{c}{Game Master} \\
    \hline

    \theutterance \stepcounter{utterance}  
    & & & \multicolumn{4}{p{0.6\linewidth}}{
        \cellcolor[rgb]{0.9,0.9,0.9}{
            \makecell[{{p{\linewidth}}}]{
                \texttt{\tiny{[P1$\langle$GM]}}
                \texttt{You are an expert in causal inference. The following question is not a typical commonsense query, but rather a meticulously designed question created by a professor specializing in causal inference, intended to assess the students' mastery of the course content.} \\
\texttt{Imagine a self{-}contained, hypothetical world with only the following conditions, and without any unmentioned factors or causal relationships: Gender has a direct effect on treatment and peanut allergy. Treatment has a direct effect on peanut allergy. For patients who are not male and not receiving treatment, the probability of being allergic to peanuts is 84\%. For patients who are not male and receiving treatment, the probability of being allergic to peanuts is 71\%. For patients who are male and not receiving treatment, the probability of being allergic to peanuts is 21\%. For patients who are male and receiving treatment, the probability of being allergic to peanuts is 3\%. The overall probability of male gender is 54\%. For patients receiving treatment, would it be more likely to see being allergic to peanuts if the patient had received no treatment?} \\
\texttt{Start your answer with 'Yes' or 'No', followed by additional reasoning or evidence to support your explanation.} \\
            }
        }
    }
     \\ \\

    \theutterance \stepcounter{utterance}  
    & \multicolumn{4}{p{0.6\linewidth}}{
        \cellcolor[rgb]{0.561,0.737,0.561}{
            \makecell[{{p{\linewidth}}}]{
                \texttt{\tiny{[P1$\rangle$GM]}}
                \texttt{No} \\
\\ 
\texttt{**Reasoning:**} \\
\\ 
\texttt{To answer this question, we must determine the causal effect of treatment on peanut allergy for the specific subgroup of patients who are currently receiving treatment. The question asks: "For patients receiving treatment, would it be more likely to see being allergic to peanuts if the patient had received no treatment?"} \\
\\ 
\texttt{In causal inference terms, this is asking for the **Counterfactual Probability**: \$P(\textbackslash{}text\{Allergy\} | \textbackslash{}text\{Treatment\}=1, \textbackslash{}text\{Gender\}=g, \textbackslash{}text\{Do\}(\textbackslash{}text\{Treatment\}=0))\$. We need to compare this counterfactual probability to the observed probability for the treated group, \$P(\textbackslash{}text\{Allergy\} | \textbackslash{}text\{Treatment\}=1, \textbackslash{}text\{Gender\}=g)\$.} \\
\\ 
\texttt{Let's break down the data provided:} \\
\texttt{1.  **Observed Probabilities (Factual World):**} \\
\texttt{    *   Non{-}male, Treated: \$P(Y=1 | X=0, T=1) = 0.71\$} \\
\texttt{    *   Male, Treated: \$P(Y=1 | X=1, T=1) = 0.03\$} \\
\\ 
\texttt{2.  **Counterfactual Probabilities (Hypothetical World where Treated patients receive No Treatment):**} \\
\texttt{    *   The problem states that "Gender has a direct effect on treatment and peanut allergy" and "Treatment has a direct effect on peanut allergy." Cruc} \\
            }
        }
    }
    & & \\ \\

    \theutterance \stepcounter{utterance}  
    & & & \multicolumn{2}{p{0.3\linewidth}}{
        \cellcolor[rgb]{0.9,0.9,0.9}{
            \makecell[{{p{\linewidth}}}]{
                \texttt{\tiny{[GM$|$GM]}}
                \texttt{No} \\
            }
        }
    }
    & & \\ \\

    \theutterance \stepcounter{utterance}  
    & & & \multicolumn{2}{p{0.3\linewidth}}{
        \cellcolor[rgb]{0.9,0.9,0.9}{
            \makecell[{{p{\linewidth}}}]{
                \texttt{\tiny{[GM$|$GM]}}
                \texttt{yes} \\
            }
        }
    }
    & & \\ \\

    \theutterance \stepcounter{utterance}  
    & & & \multicolumn{2}{p{0.3\linewidth}}{
        \cellcolor[rgb]{0.9,0.9,0.9}{
            \makecell[{{p{\linewidth}}}]{
                \texttt{\tiny{[GM$|$GM]}}
                \texttt{game\_result = LOSE} \\
            }
        }
    }
    & & \\ \\

\end{supertabular}
}

\end{document}
